\chapter{Shared-Signal Summary}
\label{app:signals}

The table below is extracted from ADR-003 and provides a compact view of the shared coordination signals. It summarizes architectural behavior only. Electrical levels, pin assignments, biasing, termination, pulse widths, timing margins, and inter-backplane implementation belong to the applicable ICDs.

% Generated from ADR-003 R3. Do not edit.
\newcommand{\GeneratedSharedSignalRows}{%
\texttt{OK} & Open-drain contributions from boards and qualified backplane rail-health sources; pull-up on main & HIGH means no participant is asserting a fault. Any participant may pull LOW. A LOW trip has priority over commanded operation. \\
\texttt{LOOP\_OUT / LOOP\_IN} & Series continuity path originating and terminating at main; passive copper or unconditional active forwarding & A broken return path causes main to assert its OK fault contribution. \\
\texttt{EN} & Main; shared level & HIGH only while the system is armed in RUN. LOW means safe/disarmed. Function boards independently qualify a rising edge; falling removes hardware permission without waiting for an FSM transition. \\
\texttt{CLEAR} & Main; shared level/pulse & Requests locally faulted boards to clear retained evidence as appropriate and re-evaluate live conditions. It does not force a board to declare itself healthy. \\
\texttt{CLOCK} & Main clock source; point-to-point LVDS & Continuous 100 MHz sequencer clock. Function boards use it directly for timing-critical work and independently detect its loss. \\
\texttt{SYNC} & Main; point-to-point LVDS & With EN=1, rising and falling acquisition events are captured in the 100 MHz domain. With EN=0, a pulse is a converter-alignment event that compatible function boards recognize in any local state. A separate CDC-safe observation is used for management. Detailed behavior is defined by ADR-004. \\
}


\begin{table}[htbp]
\centering
\footnotesize
\begin{tabularx}{\textwidth}{@{}>{\raggedright\arraybackslash}p{0.17\textwidth}>{\raggedright\arraybackslash}p{0.27\textwidth}>{\raggedright\arraybackslash}X@{}}
\toprule
\textbf{Signal} & \textbf{Driver and topology} & \textbf{Architectural behavior} \\
\midrule
\GeneratedSharedSignalRows
\bottomrule
\end{tabularx}
\caption{Shared-signal behavior extracted from ADR-003 R3 during the publication build.}
\end{table}

Five point-to-point utility-converter synchronization pairs, numbered 0 through 4, are also part of the standard main-board/backplane interface. They synchronize backplane utility-converter channels during acquisition. Their signal names, mapping, divisors, phase behavior, and electrical implementation are defined by the backplane ICD under ADR-005.

A function board may additionally derive compatible local converter clocks from its point-to-point \texttt{CLOCK}. A \texttt{SYNC} pulse received while \texttt{EN=0} may align those local dividers in any local board state. With \texttt{EN=1}, \texttt{SYNC} has acquisition meaning only. This optional board-local use is defined by ADR-004 and does not consume or redefine the five backplane utility-converter synchronization pairs.

Board-local watchdog, fail-safe, relay-control, and voltage-supervisor signals are not shared backplane signals. Their internal names and implementation belong to board design specifications.
